\chapter{Точное наблюдаемое чтение и внедиагональный юкавский угол}

Коммутирующий квадрат разделил семейное и наблюдаемое чтения, но оставил
открытой природу взаимодействия между ними. Проверим, может ли прежняя
проекция, привязанная к состоянию, быть условным ожиданием того же квадрата.

\section{Наблюдаемый пятнадцатимерный блок}

Минимальный пакет без правого нейтрино имеет разложение
\begin{equation}
 H_{15}=Q_L^{(6)}\oplus L_L^{(2)}\oplus u_R^{(3)}
 \oplus d_R^{(3)}\oplus e_R^{(1)}.
\end{equation}
Оно несёт координатную алгебру
\begin{equation}
 \mathcal A_{\mathrm{SM}}=\mathbb C\oplus\mathbb H\oplus M_3(\mathbb C)
\end{equation}
и после унимодулярности калибровочную алгебру
$\mathfrak{su}(3)\oplus\mathfrak{su}(2)\oplus\mathfrak u(1)$.
Проекция на пять бимодульных блоков
\begin{equation}
 E_{15}(T)=\sum_{s}P_sTP_s
\end{equation}
унитальна, сохраняет след, инволюцию и идемпотентна. Все машинные дефекты
равны нулю. Поскольку она действует только на наблюдаемом множителе, она
коммутирует с семейным условным ожиданием предыдущей главы.

\section{Повторный аудит проекции, привязанной к состоянию}

Ранее была построена ортогональная проекция
\begin{equation}
 \Pi_\rho(T)=T-QTQ,
 \qquad Q=I-\rho,
 \label{eq:v5-state-anchored-projection}
\end{equation}
на пятимерный бимодуль
$\rho M_3+M_3\rho$. Она идемпотентна и самосопряжённа относительно
гильбертово-шмидтовской метрики. Однако это не делает её положительной.

В базисе $\rho=\operatorname{diag}(1,0,0)$ положительная матрица
\begin{equation}
 T_+=
 \begin{pmatrix}1&1&1\\1&1&1\\1&1&1\end{pmatrix}
\end{equation}
имеет спектр $\{3,0,0\}$, тогда как
\begin{equation}
 \Spec\Pi_\rho(T_+)=\{2,0,-1\}.
\end{equation}
Матрица Чоя отображения~\eqref{eq:v5-state-anchored-projection} также имеет
собственное значение $-1$. Следовательно, $\Pi_\rho$ не положительна, не
вполне положительна и не является условным ожиданием.

\begin{theorem}[Исправление статуса проекции]
Проекция $\Pi_\rho$ является канонической ортогональной проекцией на
операторный бимодуль, но не картой состояний. Прежнее наименование её
вполне положительной картой отзывается.
\end{theorem}

Это согласуется с общей теорией: условное ожидание на $C^*$-подалгебру
является унитальной вполне положительной проекцией. Юкавский оператор не
обязан быть картой состояний и потому должен быть помещён в другой угол
родительского объекта.

\section{Что даёт настоящее условное ожидание}

Единственное сохраняющее след ожидание на алгебру состояния
$\mathbb C\rho\oplus\mathbb C Q$ имеет вид
\begin{equation}
 E_\rho(T)=\rho T\rho+\frac12\Tr(QTQ)Q.
\end{equation}
Его матрица Чоя неотрицательна, а дефекты унитальности, идемпотентности и
сохранения следа не превосходят $5\cdot10^{-16}$. Но оно уничтожает все
операторы между $\rho$ и $Q$ и поэтому не может породить юкавское смешивание.

Отсюда следует структурное разделение:
\begin{enumerate}
 \item диагональные чтения состояний задаются условными ожиданиями;
 \item взаимодействия задаются внедиагональными операторными бимодулями.
\end{enumerate}

\section{Связывающая алгебра}

Стандартным контейнером для двух алгебр $A,B$ и бимодуля
${}_AE_B$ является связывающая алгебра
\begin{equation}
 \mathcal L(E)=
 \begin{pmatrix}
 A&E\\ E^*&B
 \end{pmatrix}.
 \label{eq:v5-linking-algebra}
\end{equation}
Диагональные углы хранят семейное и наблюдаемое чтения, а внедиагональный
угол --- переходы между ними. Произведения $EE^*\subset A$ и
$E^*E\subset B$ обеспечивают замкнутость и позволяют нормировать все углы
одним следом.

В нашей программе пятимерный модуль
$\rho M_3+M_3\rho$ теперь следует читать именно как кандидат на такой
внедиагональный угол, а не как условное ожидание. Его оставшееся условие
$QTQ=0$ ещё должно быть выведено из аффинного проектора, а не принято как
аксиома.

\begin{protocol}
\begin{enumerate}
 \item точное пятнадцатимерное чтение Стандартной модели:
 \textbf{встроено};
 \item совместимость его ожидания с семейным фактором:
 \textbf{получена};
 \item статус $\Pi_\rho$ как вполне положительной карты:
 \textbf{опровергнут};
 \item статус $\Pi_\rho$ как проекции на операторный бимодуль:
 \textbf{сохранён};
 \item настоящее положительное ожидание $E_\rho$:
 \textbf{получено, но смешивание уничтожает};
 \item связывающая алгебра с юкавским внедиагональным углом:
 \textbf{выделена для следующей проверки};
 \item происхождение условия опоры из аффинного проектора:
 \textbf{открыто};
 \item физическое замыкание:
 \textbf{не достигнуто}.
\end{enumerate}
\end{protocol}

Машинная проверка и сертификат сохранены в
\path{s2t_v5_sm_linking_corner_gate.py} и
\path{s2t_v5_sm_linking_corner_gate_results.json}.